%%%%%%%%%%%%%%%%%%%%%%%%%%%%%%%%%%%%%%%%%%%%%%%%%%%%%%%%%%%%%%%%%%%%%%
% Lớp 10 | Phương pháp tọa độ trong mặt phẳng (𝑂𝑥𝑦) | Tích vô hướng (theo tọa độ)
% Dạng 0H9?2-1: Tích vô hướng, góc giữa 2 véctơ
% Số câu: 7  (sinh tự động bởi tach_cau.py; từ nay thầy sửa trực tiếp file này)
%%%%%%%%%%%%%%%%%%%%%%%%%%%%%%%%%%%%%%%%%%%%%%%%%%%%%%%%%%%%%%%%%%%%%%

%%%%%%%%%%%%%%%% TRẮC NGHIỆM 4 PHƯƠNG ÁN %%%%%%%%%%%%%%%%

% nguồn: MaHoaNH_TracNghiem_Lop10
\begin{ex}%[0H9N2-1]%[ma:df14f19d]
	Trong mặt phẳng $Oxy$,  cho $\overrightarrow{a}=(1;3)$, $\overrightarrow{b}=(-2;1)$. Khi đó $\overrightarrow{a}\cdot\overrightarrow{b}$ bằng
	\choice 
		 {$3$}
		 {$2$}
		 {$4$}
		{\True  $1$}
	 \loigiai{
		Ta có $\overrightarrow{a}\cdot\overrightarrow{b}=1\cdot(-2)+3\cdot  1=1$.
	 }
\end{ex}

% nguồn: MaHoaNH_TracNghiem_Lop10
\begin{ex}%[0H9H2-1]%[ma:622ac075]
	 Trong mặt phẳng $Oxy$, tích vô hướng của hai vectơ $\vec{a}=(1; -2)$, $\vec{b}=(-5; 6)$ bằng
	 \choice
		 {$7$}
		{\True $-17$} 
		{$-4$} 
		{$16$} 
	\loigiai{ 
		Ta có  $\vec{a}\cdot\vec{b}=1\cdot(-5)+(-2)\cdot6=-5-12=-17$.
	 }
\end{ex}


%%%%%%%%%%%%%%%% ĐÚNG/SAI %%%%%%%%%%%%%%%%

% nguồn: MaHoaNH_DungSai_Lop10
\begin{ex}%[0H9N2-1]%[ma:2373d24d]
	Trong hệ tọa độ $Oxy$, cho các  điểm $A(0;3)$, $B(4;0)$, $C(-2;-5)$. Mỗi kết quả dưới đây đúng hay sai?
	 \choiceTF
		 {\True $\overrightarrow{AB}\cdot\overrightarrow{AC}=16$}
		 {$\overrightarrow{OB}\cdot\overrightarrow{CA}=9$}
		{$\overrightarrow{CA}\cdot\overrightarrow{CB}=-10$} 
		{\True  $\overrightarrow{AB}\cdot\overrightarrow{BC}=-9$}
	\loigiai{ 
		Ta có $\overrightarrow{AB}=(4;-3)$,  $\overrightarrow{AC}=(-2;-8)$, $\overrightarrow{BC}=(-6;-5)$, $\overrightarrow{OB}=(4;0)$.
		\begin{itemchoice}
			\itemch $\overrightarrow{AB}\cdot\overrightarrow{AC}=4\cdot (-2)+(-3)\cdot (-8)=16$. 
			\itemch  $\overrightarrow{OB}\cdot\overrightarrow{CA}=4\cdot 2+0\cdot 8=8$.
			\itemch  $\overrightarrow{CA}\cdot\overrightarrow{CB}=2\cdot 6+8\cdot 5=52$.
			 \itemch $\overrightarrow{AB}\cdot\overrightarrow{BC}=4\cdot (-6)+(-3)\cdot (-5)=-9$.
		\end{itemchoice}
	 }
\end{ex}

% nguồn: MaHoaNH_DungSai_Lop10
\begin{ex}%[0H9N2-1]%[0H9H3-2]%[ma:e4b3a25e]
	Trong mặt phẳng tọa độ $Oxy$, cho  tam giác $ABC$ có $A(1;1)$, $B(0;-2)$, $C(4;2)$. Xét tính đúng, sai của các mệnh đề sau
	 \choiceTF
		{\True $\overrightarrow{AB}=(-1;-3)$} 
		{Điểm $B$ thuộc đường thẳng $\Delta\colon  2x-3y+1=0$}
		{$\overrightarrow{AB}\cdot\overrightarrow{AC}=0$} 
		{\True Phương trình đường  trung tuyến $AM$ là $x+y-2=0$}
	 \loigiai{
		Tọa độ  véc-tơ $\overrightarrow{AB}=(-1;-3)$.\\
		Thay điểm $B(0;-2)$ vào đường thẳng  $\Delta$ suy ra $B\notin\Delta$.\\
		Ta có  $\overrightarrow{AB}\cdot\overrightarrow{AC}=(-1;-3)\cdot(3;1)=-6$.\\
		Trung điểm của $BC$ là điểm $M$ có tọa độ $(2;0)$. Phương trình đường trung tuyến $AM$ là $$\dfrac{x-1}{2-1}=\dfrac{y-1}{0-1}\Leftrightarrow  x+y-2=0.$$
		\begin{itemchoice}
			\itemch $\overrightarrow{AB}=(-1;-3)$ đúng. 
			\itemch Điểm $B$ thuộc đường thẳng  $\Delta\colon 2x-3y+1=0$ sai.
			\itemch $\overrightarrow{AB}\cdot\overrightarrow{AC}=0$  sai.
			\itemch Phương trình  đường trung tuyến $AM$ là $x+y-2=0$ đúng.
		\end{itemchoice}
	 }
\end{ex}

% nguồn: MaHoaNH_DungSai_Lop10
\begin{ex}%[0H9H2-1]%[ma:9db899af]
	Trong hệ tọa độ $Oxy$, cho hai véc-tơ $\overrightarrow{a}=(1;1)$ và $\overrightarrow{b}=(2;1)$. Biết $\overrightarrow{c}$ thỏa mãn $\left(\overrightarrow{a}+\overrightarrow{b}\right)\cdot\overrightarrow{c}=-1$ và $\left(\overrightarrow{a}+2\overrightarrow{c}\right)\cdot\overrightarrow{b}=1$. Trong các  khẳng định sau, khẳng định nào đúng, khẳng định nào sai?
	\choiceTF 
		{$\left|\overrightarrow{c}\right|=\sqrt{5}$} 
		{\True $\overrightarrow{b}\cdot\overrightarrow{c}=-1$} 
		{\True  $\overrightarrow{a}\perp\overrightarrow{c}$}
		 {\True $\left(\overrightarrow{a}-\overrightarrow{b}\right)\cdot\overrightarrow{c}=1$}
	\loigiai{ 
		Giả sử $\overrightarrow{c}=(x;y)$, thay vào giả  thiết ta được
		$$\heva{& \left(\overrightarrow{a}+\overrightarrow{b}\right)\cdot\overrightarrow{c}=3x+2y\\&  \left(\overrightarrow{a}+2\overrightarrow{c}\right)\cdot\overrightarrow{b}=4x+2y+3} \Leftrightarrow \heva{& 3x+2y=-1\\& 4x+2y=-2} \Leftrightarrow \heva{& x=-1\\& y=1.}$$
		Vậy $\overrightarrow{c}=(-1;1)$. 
		\begin{itemchoice}
			 \itemch $\left|\overrightarrow{c}\right|=\sqrt{(-1)^2+1^2}=\sqrt{2}$.
			\itemch $\overrightarrow{b}\cdot\overrightarrow{c}=2\cdot (-1)+1\cdot 1=-1$. 
			\itemch $\overrightarrow{a}\cdot\overrightarrow{c}=1\cdot (-1)+1\cdot  1=0 \Rightarrow \overrightarrow{a}\perp\overrightarrow{c}$.
			\itemch  Ta có $\overrightarrow{a}-\overrightarrow{b}=(-1;0)$, suy ra $\left(\overrightarrow{a}-\overrightarrow{b}\right)\cdot\overrightarrow{c}=(-1)\cdot (-1)+1\cdot 0=1$.
		\end{itemchoice}
	} 
\end{ex}

% nguồn: MaHoaNH_DungSai_Lop10
\begin{ex}%[0H9V2-1]%[ma:7c3a6466]
	Trong mặt phẳng tọa độ $Oxy$, cho ba véc-tơ  $\overrightarrow{a}=(3;-1)$, $\overrightarrow{b}=(1;-2)$, $\overrightarrow{c}=(-1;7)$.
	\choiceTF 
		{$\overrightarrow{a}\perp  \overrightarrow{b}$}
		{\True Tọa độ véc-tơ $\overrightarrow{u}=\overrightarrow{a}-2\overrightarrow{b}+3\overrightarrow{c}$ là  $\overrightarrow{u}=(-2;24)$}
		{\True Cô-sin của góc  giữa hai véc-tơ $\overrightarrow{a}$ và $\overrightarrow{b}$ bằng $\dfrac{1}{\sqrt{2}}$}
		{Biết rằng $\overrightarrow{c}=m\overrightarrow{a}+n\overrightarrow{b}$ (với  $m,n\in \mathbb{R}$). Khi đó $m-n=-3$}
	\loigiai{ 
		\begin{itemchoice}
			\itemch Vì $3\cdot 1+(-1)\cdot (-2)=5\ne 0$ nên $\vec{a}$ không  vuông góc $\vec{b}$.
			 \itemch Ta có $2\overrightarrow{b}=(2;-4)$, $3\overrightarrow{c}=(-3;21)$.\\
			Vậy $\overrightarrow{a}-2\overrightarrow{b}+3\overrightarrow{c}=(-2;24)$. 
			\itemch Ta có 
			$$\cos\left(\vec{a};\vec{b}\right)=\dfrac{|\vec{a}\cdot  \vec{b}|}{|\vec{a}|\cdot |\vec{b}|}=\dfrac{5}{\sqrt{3^2+(-1)^2}\cdot \sqrt{1^2+(-2)^2}}=\dfrac{1}{\sqrt{2}}.$$
			\itemch Ta có $\overrightarrow{c}=m\overrightarrow{a}+n\overrightarrow{b} \Leftrightarrow \heva{&3m+n=-1\\&-m-2n=7} \Leftrightarrow  \heva{&m=1\\&n=-4}$\\
			 Vậy $m-n=5$.
		\end{itemchoice}
	 }
\end{ex}


%%%%%%%%%%%%%%%% TRẢ LỜI NGẮN %%%%%%%%%%%%%%%%

% nguồn: MaHoaNH_TraLoiNgan_Lop10
\begin{bt}%[0H9H2-1]%[ma:f7a4fad0]
	Trong mặt phẳng tọa độ $Oxy$ cho hai  điểm $A(3;6)$, $B(x;-2)$ với $x\in\mathbb{R}$. Tìm $x$ để $OA$ vuông góc với $AB$.
	\shortans{19} 
	\loigiai{ 
		Ta có  $\overrightarrow{OA}=(3;6)$ và $\overrightarrow{AB}=(x-3;-8)$.\\
		Để $\overrightarrow{OA}\perp \overrightarrow{AB}$ thì $3(x-3)-48=0\Leftrightarrow x=19$. 
	} 
\end{bt}
